% Generated by alp.report -- regenerate rather than edit by hand.
% Data provenance: SIMULATED.
% WARNING: reference-model output. Replace with campaign results before submission.
\subsection{Experimental Setup}
The field evaluation was conducted on the two-site cyber range of the University of Customs and Finance, with the sites interconnected through a protected VPN. 25 Kali Linux workstations generated identical precomputed transaction traces. A permissioned Avalanche L1 based on Subnet-EVM used five validators and two independent read nodes. We compared the stock configuration with minimum block-delay profiles of 1000, 750, 500 and 250\,ms under local, two-site VPN, and emulated three-region topologies. Each condition was repeated 3 times using a randomized blocked schedule. A run comprised a 5\,s warm-up, a 20\,s measurement window and a 5\,s drain period. The primary metric, $T_\mathrm{visible}$, was measured on the load generator using a monotonic nanosecond clock from raw-transaction submission until the first independent read returned the accepted updated state. Run-level $p_{50}$, $p_{95}$ and $p_{99}$ values were compared using paired bootstrap confidence intervals.

\subsection{Results}
On the T0 local topology the best static configuration was C4, stable up to 400\,tx/s with a mean $p_{99}$ of 278\,ms.
On the T1 physical VPN topology the best static configuration was C4, stable up to 400\,tx/s with a mean $p_{99}$ of 345\,ms.
On the T2 three-region (emulated) topology the best static configuration was C3, stable up to 400\,tx/s with a mean $p_{99}$ of 793\,ms.
Hypothesis H1 was supported: C4: 250 ms; C3: 500 ms; C2: 750 ms; C1: 1000 ms.
Hypothesis H2 was not supported: realisation efficiency 95 \% at 25 tx/s vs 96 \% at 200 tx/s (predicted gain 742 ms from C1 to C4).
Hypothesis H3 was not supported: C4 kept the lowest p99 in every stratum.
